\documentclass[11pt]{article}
\usepackage[T1]{fontenc}
\usepackage[utf8]{inputenc}
\usepackage{lmodern,microtype}
\usepackage[margin=1in,headheight=15pt]{geometry}
\usepackage{amsmath,amssymb,amsthm,mathtools}
\usepackage{booktabs,tabularx,array,longtable}
\usepackage{enumitem,xcolor,listings,fancyhdr}
\usepackage{hyperref}
\definecolor{accent}{HTML}{455846}
\definecolor{codebg}{HTML}{F5F5F1}
\hypersetup{colorlinks=true,linkcolor=accent,citecolor=accent,urlcolor=accent,
 pdftitle={From algebraic roots to radicals},pdfauthor={Technical report prepared for Vladimir Reshetnikov}}
\lstset{basicstyle=\ttfamily\small,backgroundcolor=\color{codebg},
 columns=fullflexible,keepspaces=true,breaklines=true,showstringspaces=false,
 frame=single,rulecolor=\color{black!15},framerule=.3pt,
 xleftmargin=5pt,xrightmargin=5pt,aboveskip=10pt,belowskip=10pt}
\pagestyle{fancy}\fancyhf{}
\fancyhead[L]{\small\sffamily RadicalSolve: theory and implementation}
\fancyhead[R]{\small\sffamily Version 0.1.0}
\fancyfoot[C]{\thepage}
\setlength{\parindent}{0pt}\setlength{\parskip}{5pt}
\setlist{nosep,leftmargin=1.6em}
\newtheorem{theorem}{Theorem}[section]
\newtheorem{proposition}[theorem]{Proposition}
\newtheorem{lemma}[theorem]{Lemma}
\newtheorem{corollary}[theorem]{Corollary}
\theoremstyle{definition}\newtheorem{definition}[theorem]{Definition}
\theoremstyle{remark}\newtheorem{remark}[theorem]{Remark}
\newcommand{\Q}{\mathbb Q}\newcommand{\C}{\mathbb C}\newcommand{\Z}{\mathbb Z}
\DeclareMathOperator{\Gal}{Gal}\DeclareMathOperator{\Res}{Res}
\DeclareMathOperator{\Tr}{Tr}\DeclareMathOperator{\Norm}{Norm}
\DeclareMathOperator{\Aut}{Aut}\DeclareMathOperator{\Span}{span}
\newcommand{\code}[1]{\texttt{#1}}
\newcommand{\zetaP}{\zeta_p}
\newcommand{\fixed}[1]{L^{#1}}
\title{\vspace{-1cm}\textbf{From algebraic roots to radicals}\\[8pt]
\large A branch-aware solver based on structural reductions\\
and constructive Galois theory}
\author{Technical report prepared for Vladimir Reshetnikov}
\date{8 September 2026 \quad\textbullet\quad Reference implementation 0.1.0}
\begin{document}
\maketitle
\begin{abstract}
We construct a radical-conversion system for exact algebraic numbers represented
by polynomial roots. The starting point is the supplied \code{FindExtension.nb}
notebook and its searches for coefficient fields of quadratic and cubic factors.
The replacement has two layers. Exact structural reductions produce short formulas
for power compositions, Dickson polynomials, and generalized reciprocal
polynomials; in particular, they solve both examples in the original question.
A general reference engine constructs a splitting field, recovers its actual
automorphisms, and performs Fourier resolvent descent along a prime-index
subnormal series. Roots of unity are handled in a formal cyclotomic tensor
algebra, so no linear-disjointness assumption is needed. A norm polynomial
supplies the hypothesis needed for branch comparisons based on root separation.
We prove finite completeness of the mathematical algorithm, explain the boundary
between that theorem and practical resource limits, and document executable
Python code and a Wolfram Language frontend with native structural methods.
The Python implementation is tested; the supplied Wolfram code has not been
executed in a Wolfram kernel in this environment.
\end{abstract}
\vspace{4pt}
\noindent\textbf{Scope.} Parameter-free algebraic numbers over $\Q$, with complex
radicals allowed. No claim of minimal radical depth, shortest formula, or
practical completion for every large solvable input is made.
\newpage
\tableofcontents
\newpage

\section{The problem and the contract}

A \code{Root} expression is an excellent exact representation of an algebraic
number, but it deliberately hides how that number was obtained. The conversion
problem is to recover a finite expression involving rational arithmetic and
rational powers, with the \emph{same complex embedding} as the input. The official
\code{ToRadicals} documentation explicitly covers all degrees through four and
notes that some higher-degree radical-solvable cases are missed
\cite{wolfram-toradicals}. Failure of that transformation is therefore not a
Galois-theoretic impossibility result.

The motivating question supplies, among other examples,
\begin{align}
 r_6&=\code{Root}[-1-\#^2-\#^3+\#^4+\#^6\mathbin{\&},2],\label{eq:r6}\\
 r_5&=\code{Root}[6+25\#-25\#^3+5\#^5\mathbin{\&},5].\label{eq:r5}
\end{align}
The examples and notebook originate with Reshetnikov \cite{question,notebook}.
The package treats this as an algorithm-construction task, not as a request to
find a convenient expression for just one root.

\subsection{Three different questions}

For an input $\alpha$, three tasks must remain separate:
\emph{decide} whether radicals exist, \emph{construct} such radicals, and
\emph{identify} the requested conjugate. A correct abstract Galois group can
settle the first question while providing no formulas for the second. An
annihilating polynomial can settle neither the second nor the third by itself:
$\sqrt2$ and $-\sqrt2$ have the same minimal polynomial.

Let $f$ be the irreducible minimal polynomial of $\alpha$ over $\Q$, and let $L$
be its splitting field. The classical criterion is
\begin{equation}\label{eq:criterion}
 \alpha\text{ is expressible by radicals over }\Q
 \quad\Longleftrightarrow\quad \Gal(L/\Q)\text{ is solvable}.
\end{equation}
This is a statement about the \emph{minimal polynomial of the selected number},
not an arbitrary reducible polynomial that happens to vanish there. The
classical cyclic-extension and solvability theorems are described in
Milne's account \cite{milne}; the explicit finite construction below gives the
sufficiency direction in the form needed by this implementation.

\subsection{Allowed output and input}

The accepted output grammar is generated by rational constants, $i$, addition,
multiplication, and $b^q$ for $q\in\Q$, interpreted using principal complex
powers. Integer negative powers include division by nonzero expressions.
Explicit undefined values are rejected. In particular,
\begin{equation}
 \zeta_n=(-1)^{2/n}
\end{equation}
is an allowed radical, whereas an unevaluated trigonometric function or
\code{Root} is not. Intermediate \code{Root}/\code{CRootOf} objects are allowed
inside the algorithm; none may escape into its radical result.

The Python interface takes a rational polynomial and a zero-based SymPy root
index. The Wolfram interface takes an exact algebraic expression, obtains its
minimal polynomial, and preserves the original embedding by a final exact
comparison. Approximate coefficients and free parameters are rejected. A
symbolic parameter introduces specialization and branch-region questions that
are outside this version's contract. An algebraic number whose defining
expression happens to use algebraic coefficients is still in scope when its
rational minimal polynomial can be obtained.

\subsection{What ``complete'' means here}

The general algorithm is not a bounded search over guessed radicals. It is a
finite construction using polynomial factorization, resultants, finite groups,
and exact algebraic embedding comparisons. Removing degree and time limits
removes the implementation's deliberate cutoffs. The completeness theorem
assumes these standard exact primitives return correctly; it is not a machine-
checked proof of the Python program or a guarantee about the running time of a
particular CAS backend. This distinction matters particularly for splitting
fields and reconstruction of minimal polynomials from large radical formulas.

The delivered Wolfram file contains the structural algorithms natively and
calls the Python general engine when required. Thus the \emph{general theorem}
applies to an implemented engine, not to a placeholder for a future Galois
routine. The Wolfram kernel execution status is stated separately in
Section~\ref{sec:tests}.

\section{What the supplied notebook contributes}

The notebook was inspected as data, without evaluating its cells. A readable,
best-effort box extraction is included beside the unchanged notebook. That
extraction is not a runnable Wolfram program and is not a replacement for the
original source.

\subsection{Factor coefficients are resolvents in disguise}

The notebook's \code{findExtension} imposes divisibility by
\begin{equation}\label{eq:quadraticfactor}
 x^2+p x+q.
\end{equation}
Equating the two coefficients of the remainder to zero and eliminating $q$
produces candidate values of $p$. If the quadratic's roots are $\alpha_i$ and
$\alpha_j$, then
\begin{equation}
 p=-(\alpha_i+\alpha_j),\qquad q=\alpha_i\alpha_j.
\end{equation}
Thus its coefficient search is already an implicit search among symmetric
functions of subsets of conjugates. The analogous cubic search uses elementary
symmetric functions of triples. This is a natural practical way to discover a
smaller field in which the original polynomial factors.

What it does not provide is a completeness argument. A useful fixed field need
not be generated by the particular coefficient selected for elimination.
Restricting factor degrees to two and three does not cover a cyclic quotient of
prime order five or seven. Even a suitable factorization does not by itself
make all coefficient-field elements radical expressions over $\Q$.

The filter involving \code{ExactNumberQ} discards explicitly rational or
rational-complex numeric coefficients; that predicate is not a general test
for algebraicity \cite{wolfram-exact}. Discarding a rational trace can also
conceal a useful extension generated by a different factor coefficient.
The new system uses exact field coordinates rather than the printed form of a
candidate coefficient as its structural invariant.

\subsection{Projections and experimental recognition}

The notebook's \code{project} looks at real and imaginary parts of conjugates.
Those can indeed reveal smaller fields. The \code{fuzz} routine takes random
subsets of roots, approximates their sums, applies \code{RootApproximant}, and
then tries \code{PossibleZeroQ}. This is useful discovery machinery, but an
integer-relation or algebraic-recognition guess is not a proof that it equals
the original sum. It must be followed by a separate exact check.

The notebook also contains repeated interactive factor/solve experiments,
large polynomials associated with inverse-beta examples, references to global
state and earlier outputs, and iterative simplification. Such a research
notebook should not be repackaged by simply exporting all definitions and
calling the result a complete solver.

\subsection{Engineering changes}

The replacement makes the trusted conditions explicit. Polynomial identities
are checked before structural reconstruction; the normal closure supplies the
actual automorphisms used in general descent; and a radical candidate is never
selected merely because its decimal approximation is close. Empty searches,
invalid input, backend failures, resource cutoffs, and proved nonsolvability
have different results. Public code does not depend on the private
\code{Simplify`} context or on notebook output history.

The notebook's \code{radicalDepth} also counts trigonometric forms for its
complexity heuristic. That is a useful aesthetic measure, but it should not be
confused with the literal radical grammar required here. The package supplies
an explicit \code{RadicalExpressionQ} predicate instead.

\section{Structural reductions before Galois closure}\label{sec:structural}

The general construction can be costly even when a short answer exists.
Consequently the automatic route first searches for exact identities of a few
cheap forms. These searches are not themselves complete; the general engine is
the fallback.

\subsection{Translation and power composition}

Normalize $f$ to be monic of degree $n$ and remove its next-to-leading term:
\begin{equation}
 c=-\frac{[x^{n-1}]f}{n},\qquad q(y)=f(y+c).
\end{equation}
Let $g$ be the gcd of the positive exponents having nonzero coefficients in
$q$. If $g>1$, write
\begin{equation}
 q(y)=P(y^g).
\end{equation}
After recursively expressing the roots $v$ of $P$, every root of $f$ has the
form
\begin{equation}\label{eq:powercomp}
 x=c+\zeta_g^j v^{1/g},\qquad 0\le j<g.
\end{equation}
The recursion lowers polynomial degree. The identity proves annihilation of
all reconstructed candidates; their requested embedding is a separate choice.
This covers, for example, shifted binomials such as $(x+1)^5+2$.

\subsection{Dickson polynomials and correlated branches}

Define
\begin{equation}
 D_0(y,a)=2,\quad D_1(y,a)=y,\quad
 D_n(y,a)=yD_{n-1}(y,a)-aD_{n-2}(y,a).
\end{equation}
Direct induction gives
\begin{equation}\label{eq:dickson}
 D_n(u+a/u,a)=u^n+(a/u)^n.
\end{equation}
For a centered polynomial, the coefficient of $y^{n-2}$ suggests
$a=-[y^{n-2}]q/n$ when $n>2$. The implementation constructs $D_n$ and checks
whether $D_n-q$ is constant; it does not accept the coefficient test alone.

If $q(y)=D_n(y,a)-b$ and $a\ne0$, put
\begin{equation}\label{eq:dicksont}
 t=\left(\frac{b+\sqrt{b^2-4a^n}}2\right)^{1/n}.
\end{equation}
Then all roots are
\begin{equation}\label{eq:dicksonroots}
 y_j=\zeta_n^j t+\frac{a}{\zeta_n^j t}.
\end{equation}
The selected radicand is nonzero: the two roots of
$T^2-bT+a^n$ have nonzero product. Formula~\eqref{eq:dicksonroots} therefore has
no hidden division by zero. When $a=0$, the code uses the binomial formula
instead.

The reciprocal term is essential. Taking independent principal $n$th roots of
the two roots of $T^2-bT+a^n$ can destroy the condition that their product is
$a$. Keeping the second term as $a/u$ carries this relation through the entire
calculation. This is the simplest example of why branch correlation belongs to
the algebraic representation, not to a final numerical cleanup step.

\subsection{Generalized reciprocal composition}

For a centered monic polynomial $q$ of degree $2m$, search for
\begin{equation}\label{eq:reciprocal}
 q(x)=x^m P(x+a/x),\qquad a\in\Q\setminus\{0\},\quad\deg P=m.
\end{equation}
The constant coefficient forces $a^m=q(0)$. Rational $m$th roots of a rational
number can be found exactly by testing numerator and denominator for perfect
powers. There are at most two candidates for $a$.

For each candidate, the coefficient of $x^{2m}$ determines the leading
coefficient of $P$, the coefficient of $x^{2m-1}$ determines the next one, and
so on. More explicitly, subtract the appropriate multiples of
$x^m(x+a/x)^j$ for $j=m,m-1,\ldots,0$. The expression is a polynomial because
$j\le m$. A zero final remainder proves~\eqref{eq:reciprocal}.

Once $y$ is a root of $P$, recover
\begin{equation}\label{eq:reciproots}
 x=\frac{y\pm\sqrt{y^2-4a}}2.
\end{equation}
This is a rational-function decomposition rather than an ordinary polynomial
composition. Searching only for $f=g\circ h$ with polynomial $h$ would miss it.
The current recognizer searches rational $a$ after the standard centering
translation; it does not claim to find every possible algebraic or M\"obius
change of variable.

\section{The two original examples, explicitly}

\subsection{The sextic: a cubic followed by a quadratic}

The decisive identity is
\begin{equation}\label{eq:sexticidentity}
 x^6+x^4-x^3-x^2-1
 =x^3\left[(x-1/x)^3+4(x-1/x)-1\right].
\end{equation}
It is the reciprocal case with $a=-1$ and
\begin{equation}
 P(y)=y^3+4y-1.
\end{equation}
Put
\begin{equation}\label{eq:sexticu}
 u=\left(\frac{9+\sqrt{849}}{18}\right)^{1/3},
 \qquad y=u-\frac4{3u}.
\end{equation}
Here $u>0$ is the ordinary positive cube root. Its cube satisfies
$t^2-t-64/27=0$, and hence direct expansion gives $y^3+4y=1$.
The answer for~\eqref{eq:r6} is
\begin{equation}\boxed{\displaystyle
 r_6=\frac12\left(u-\frac4{3u}
       +\sqrt{\left(u-\frac4{3u}\right)^2+4}\right).}\label{eq:sexticanswer}
\end{equation}
Numerically this is approximately $1.1306854454620405872$.

There is no ambiguous real-root choice in this derivation. The cubic is
strictly increasing and has a unique real root. The function $x-1/x$ is
strictly increasing on each real half-line and maps each half-line onto all
of $\mathbb R$. Thus the sextic has exactly two real roots; the plus sign in
\eqref{eq:sexticanswer} gives the positive one requested in the question.
The two real roots have product $-1$.

A transparent Wolfram form is
\begin{lstlisting}
With[{u = ((9 + Sqrt[849])/18)^(1/3)},
  (u - 4/(3 u) + Sqrt[(u - 4/(3 u))^2 + 4])/2
]
\end{lstlisting}
The formula was independently checked in the Python environment by minimal-
polynomial reconstruction and embedding comparison. The same Wolfram expression
is supplied as an example, without claiming a Wolfram-kernel execution test.

\subsection{The quintic: a Dickson equation}

The normalized quintic is
\begin{equation}
 x^5-5x^3+5x+\frac65=D_5(x,1)+\frac65.
\end{equation}
Take
\begin{equation}
 v=\left(\frac{-3+4i}{5}\right)^{1/5}
\end{equation}
with its principal power. The selected root is
\begin{equation}\boxed{\displaystyle r_5=v+v^{-1}.}\label{eq:quinticanswer}
\end{equation}
Indeed, $v^5+v^{-5}=-6/5$ and~\eqref{eq:dickson} proves the equation.
Since $|(-3+4i)/5|=1$, the inverse term is the conjugate of $v$; this is also
the two-conjugate-radical expression in the original question. The result is
approximately $1.8070599950854245021$.

If $\theta=\arg(-3+4i)\in(\pi/2,\pi)$, the five reconstructed roots are
$2\cos((\theta+2\pi j)/5)$. This observation identifies the principal choice
as the largest real root. It is only an explanation of the embedding: the
returned expression~\eqref{eq:quinticanswer} contains no trigonometric function.
The implementation enumerates all five correlated branches and checks the
selected embedding.

\section{Constructing the exact normal closure}\label{sec:closure}

The remaining sections describe the general engine. Let $f\in\Q[x]$ be monic,
irreducible, and of degree $n>1$. We work with fields
\begin{equation}
 K=\Q[t]/(h(t)),\qquad h\text{ monic and irreducible}.
\end{equation}
Elements are rational coefficient vectors of length at most $d=\deg h$;
addition and multiplication are reduced modulo $h$. The initial field is a
root field of $f$.

\subsection{Adjoin a root of a remaining factor}

Factor $f$ over $K$. If every factor is linear, $K$ is a splitting field of
$f$, because it was built using roots of $f$ and contains all of them. Otherwise
choose an irreducible factor
\begin{equation}
 g(Y)=Y^e+g_{e-1}(t)Y^{e-1}+\cdots+g_0(t),\qquad e>1.
\end{equation}
Adjoining a root $\beta$ gives a field of degree $de$ over $\Q$.
To keep using a single-generator representation, try
\begin{equation}
 T=t+c\beta,\qquad c=1,2,3,\ldots,
\end{equation}
and compute
\begin{equation}\label{eq:primitive-resultant}
 H_c(T)=\Res_u\left(h(u),\ c^e g((T-u)/c,u)\right).
\end{equation}
The notation $g(Y,u)$ means that its coefficient vectors are read as
polynomials in $u$.

\begin{lemma}\label{lem:primitive}
If $H_c$ is square-free of degree $de$, then $T=t+c\beta$ is primitive for
$K(\beta)/\Q$, and the monic polynomial $H_c$ is its minimal polynomial.
A suitable positive integer $c$ exists.
\end{lemma}
\begin{proof}
There are $de$ embeddings of the separable field $K(\beta)$ into $\C$.
The resultant is the product of $T-\tau(t)-c\tau(\beta)$ over those embeddings.
It is monic of degree $de$. Square-freeness says that these $de$ values are
distinct, so $t+c\beta$ has $de$ conjugates and generates the whole field.

For two distinct embeddings $\tau,\rho$, the equality
$\tau(t)+c\tau(\beta)=\rho(t)+c\rho(\beta)$ either has no solution in $c$ or
has exactly one. It cannot hold identically, because $t$ and $\beta$ generate
the field. There are finitely many pairs. Avoiding their exceptional values
leaves a positive integer $c$.
\end{proof}

This test prevents a common primitive-element mistake: accepting a resultant
of the expected degree without checking that repeated conjugates have not
collapsed the actual generated field. The number of exceptional values is at
most $\binom{de}{2}$, giving a finite bound for the search in principle.

\subsection{Why embedding changes are harmless at this stage}

After computing $H_c$, the implementation constructs the next field with a
chosen exact root of $H_c$, then factors the original $f$ again. It need not
preserve the numerical embedding of every previous generator. Any root of the
proved irreducible $H_c$ defines an isomorphic field, and the polynomial $f$
still has rational coefficients. Splitting $f$ is an embedding-independent
property.

This is deliberately different from repeatedly calling a generic primitive-
element routine on a large list of already embedded \code{CRootOf} expressions.
Here the minimal polynomial has already been proved and is passed explicitly
to the field constructor. It avoids asking the backend to rediscover it.
Only after the full field is available is a concrete embedding used for radical
branch selection and for matching the original requested root.

\subsection{Termination of the closure construction}

Every adjunction is generated by a root of $f$, so the resulting abstract field
embeds into a splitting field of $f$. Each nonlinear factor adjunction increases
the degree by at least a factor of two. The degree is bounded by $n!$, and after
at most $n-1$ distinct roots have been adjoined, the remaining root is determined
by the trace coefficient. Consequently the construction reaches a splitting
field in finitely many steps.

This gives a degree bound, not a useful worst-case speed guarantee. The
coefficient heights of $H_c$ can grow substantially. Exact factorization over a
number field may dominate even when the ultimate Galois group has modest order.
The implementation checks its field-degree budget before making each extension.

\subsection{A worked closure: \texorpdfstring{$x^3-2$}{x cubed minus two}}

Start with $t^3=2$. Over $\Q(t)$,
\begin{equation}
 x^3-2=(x-t)(x^2+tx+t^2).
\end{equation}
For the quadratic factor, the first candidate gives
\begin{equation}
 H_1(T)=(T^3+2)^2,
\end{equation}
which correctly fails the square-free test. The second gives
\begin{equation}
 H_2(T)=T^6+108.
\end{equation}
The degree-six field splits $x^3-2$. The repeated polynomial at $c=1$ has a
simple interpretation: the sum of two roots is the negative of the third,
so it generates only a cubic field. The weighted sum with coefficients $1,2$
separates all six embeddings.

\section{Recovering the actual automorphisms}\label{sec:autos}

Knowing that the Galois group is isomorphic to, for example, $S_3$ is not enough
for the next stage. We need to apply each automorphism to an element represented
in the chosen field basis.

\subsection{Weighted conjugate sums instead of a second hard factorization}

One option is to factor the primitive polynomial $h$ over $L$. That is correct,
but factoring a degree-$d$ polynomial over a degree-$d$ field can be unnecessarily
expensive. The implementation uses information already accumulated during
closure construction.

The final primitive element has the form
\begin{equation}\label{eq:weighted}
 t=w_1\alpha_{i_1}+\cdots+w_r\alpha_{i_r},
 \qquad w_1=1,\quad w_j\in\Z_{>0},
\end{equation}
for distinct roots of the original $f$. The weights are the primitive-search
coefficients chosen at successive steps. Although its concrete embedding may
have changed, the existence of such a representation is preserved.

Enumerate ordered tuples of $r$ distinct original roots, form the corresponding
weighted sums $b\in L$, and retain precisely those satisfying $h(b)=0$ in
$\Q[t]/(h)$. Deduplicate by exact coefficient vectors. Stop when $d$ distinct
values have been found.

\begin{proposition}\label{prop:autos}
The procedure recovers exactly the $d$ images of $t$ under $\Gal(L/\Q)$.
Each retained $b$ defines an automorphism by $A(t)\mapsto A(b)$ modulo $h$.
\end{proposition}
\begin{proof}
Every embedding of $L$ permutes the roots of $f$ and sends~\eqref{eq:weighted}
to one of the enumerated sums. Thus all conjugates of $t$ occur. Conversely,
$h(b)=0$ gives an injective $\Q$-homomorphism $\Q[t]/(h)\to L$. Its image has
degree $d$, so it is all of $L$. There are exactly $d$ distinct conjugates,
which is both a completeness check and a stopping criterion.
\end{proof}

This turns a feature of the notebook's conjugate-sum experiments into a finite,
exact construction: the weights are known, the search space is explicit, and
the acceptance test is a zero coefficient vector, not algebraic recognition
from a decimal approximation.

\subsection{The multiplication table}

If $b_a,b_b$ are the primitive images for $\sigma_a,\sigma_b$, then
\begin{equation}
 (\sigma_a\circ\sigma_b)(t)=b_b(b_a(t))\pmod {h(t)}.
\end{equation}
The apparent reversed order is important: the outer automorphism acts on the
coefficient representation of the inner image. Horner substitution constructs
the product table and inverse table. The identity is the element whose image
is the residue class $t$ itself.

For a field element $A(t)$ and a subgroup $H$, membership in its fixed field is
checked directly:
\begin{equation}\label{eq:fixedtest}
 A\in L^H\quad\Longleftrightarrow\quad
 A(b_\sigma(t))\equiv A(t)\pmod h\quad(\sigma\in H).
\end{equation}
No smaller fixed-field primitive polynomial is required to carry out this test.
The representation may be redundant, but it is uniform and exact.

\section{From a solvable group to prime cyclic steps}\label{sec:chain}

The descent needs a subnormal chain
\begin{equation}\label{eq:chain}
 G=H_0\triangleright H_1\triangleright\cdots\triangleright H_s=1,
 \qquad H_{i-1}/H_i\cong C_{p_i},
\end{equation}
where $p_i$ is prime. These subgroups need only be normal in their immediate
predecessor, not in the whole group.

For a current subgroup $H$, compute its commutator subgroup $D=[H,H]$ by finite
group closure. If $D=H\ne1$, a nontrivial perfect subgroup has been found and
$G$ is not solvable. Otherwise begin with $N=D$, scan the elements $g\in H$,
and replace $N$ by $\langle N,g\rangle$ whenever that generated subgroup is
still proper in $H$.

\begin{lemma}
At the end of this scan, $N$ is normal in $H$ and $H/N$ has prime order.
\end{lemma}
\begin{proof}
The quotient $H/D$ is abelian, so every subgroup containing $D$ is normal in
$H$. If a tested $g$ generated all of $H$ together with an earlier $N$, it still
does so after $N$ grows. Thus at the end every $g\notin N$ generates the whole
quotient $H/N$. That finite quotient has no nontrivial proper subgroup, and
is therefore cyclic of prime order.
\end{proof}

Repeat with $H=N$. The code explicitly checks normality and prime index before
using a step. Any $\sigma\in H_{i-1}\setminus H_i$ generates the quotient,
and $\sigma^{p_i}$ fixes every element of $L^{H_i}$.

For degrees through six, the implementation optionally asks SymPy's exact
Galois classifier for an early negative result. Its documented degree limit
is six \cite{sympy-fields}. That classifier is \emph{not} the source of the
actual automorphism action used for positive construction. Above degree six,
the explicit closure and multiplication-table algorithm still supplies that
action. Exhaustion of the classifier's transformation attempts falls back to
the general route instead of implying nonsolvability.

\section{Fourier descent without a disjointness assumption}\label{sec:fourier}

Fix one step $N\triangleleft H$ with $H/N\cong C_p$, and let
$\beta\in L^N$. Choose $\sigma\in H$ generating the quotient. Classical
Lagrange resolvents diagonalize the cyclic action after a primitive $p$th root
of unity is available. This is the computational content behind the cyclic
Kummer construction \cite{milne}. Related general radical-solution algorithms
are discussed by Li \cite{li}; the presentation here focuses on exact field
coordinates, a division-free tensor formulation, and explicit branch bookkeeping.

\subsection{The tensor algebra}

Work formally in
\begin{equation}\label{eq:tensor}
 A=L[z]/(\Phi_p(z))
   \cong L\otimes_{\Q}\Q(\zeta_p),\qquad
 \Phi_p(z)=1+z+\cdots+z^{p-1}.
\end{equation}
This is a free $L$-module with basis $1,z,\ldots,z^{p-2}$. It is \emph{not
necessarily a field}. For example, when $L$ already contains $\zeta_3$,
\begin{equation}
 (z-\zeta_3)(z-\zeta_3^2)=0
\end{equation}
in $A$, although neither factor is zero. Treating such a quotient as a field
and inverting arbitrary nonzero elements would be a real correctness bug.

The implementation uses only addition, multiplication, and positive integer
powers in $A$. Its coefficients remain elements of the exact field $L$.
Reduction uses
\begin{equation}
 z^{p-1}=-(1+z+\cdots+z^{p-2}),\qquad z^p=1.
\end{equation}
No tensor-field inversion or linear-disjointness test is needed.

\subsection{The resolvent identities}

For $0\le k<p$, define
\begin{equation}\label{eq:resolvent}
 R_k(z)=\sum_{j=0}^{p-1}z^{-kj}\sigma^j(\beta).
\end{equation}
Negative powers can be reduced modulo $p$ because $z^p=1$. The action of $H$
on $A$ fixes the formal symbol $z$ and acts on coefficients.

\begin{proposition}\label{prop:descent}
The following identities hold in the tensor algebra:
\begin{align}
 R_0&\in L^H,\label{eq:trace}\\
 \sigma R_k&=z^kR_k,\label{eq:eigen}\\
 R_k^p&\in L^H[z]/(\Phi_p(z))\quad(1\le k<p),\label{eq:powerfixed}\\
 p\beta&=\sum_{k=0}^{p-1}R_k.\label{eq:inverseDFT}
\end{align}
\end{proposition}
\begin{proof}
Normality of $N$ implies that every conjugate $\sigma^j(\beta)$ is fixed by
$N$. The relation $\sigma^p(\beta)=\beta$ makes a cyclic index shift in
\eqref{eq:resolvent} valid and gives~\eqref{eq:eigen}. Its $p$th power is
fixed by $\sigma$ and also by $N$, hence by $H$. Since the displayed basis of
$A$ is free over $L$, invariance of an element is equivalent to invariance of
each coefficient; this proves~\eqref{eq:powerfixed} even if $A$ has zero divisors.
The same argument for $k=0$ proves~\eqref{eq:trace}.

Finally, summing~\eqref{eq:resolvent} over $k$ gives $p\beta$ because
$\sum_{k=0}^{p-1}z^{-kj}=0$ for $1\le j<p$, and equals $p$ for $j=0$.
For nonzero $j$, primality of $p$ permutes the nontrivial exponents modulo $p$,
so this vanishing is exactly the relation $\Phi_p(z)=0$.
\end{proof}

\subsection{Recursive conversion}

Write
\begin{equation}\label{eq:qcoeffs}
 R_k(z)^p=\sum_{\ell=0}^{p-2}q_{k\ell}z^\ell,
 \qquad q_{k\ell}\in L^H.
\end{equation}
The trace and these coefficients belong to the preceding fixed-field level
and can be recursively converted. Once their radical expressions are known,
evaluate $z=\zeta_p$ and form
\begin{equation}
 Q_k=\sum_{\ell=0}^{p-2}\operatorname{Rad}(q_{k\ell})\zeta_p^\ell.
\end{equation}
There is an integer $e_k\in\{0,\ldots,p-1\}$ for which
\begin{equation}\label{eq:branchresolvent}
 R_k(\zeta_p)=\zeta_p^{e_k}Q_k^{1/p}.
\end{equation}
The next section explains how that integer is selected without numerical
guessing. The reconstructed expression is
\begin{equation}\label{eq:radicalrecursion}
 \operatorname{Rad}(\beta)
 =\frac1p\left[\operatorname{Rad}(R_0)+
      \sum_{k=1}^{p-1}\zeta_p^{e_k}Q_k^{1/p}\right].
\end{equation}
At the bottom of the subgroup chain, the fixed elements are rational coefficient
vectors. The recursion therefore terminates. Values already fixed by the
preceding, larger subgroup skip a level, and equal coefficient vectors at the
same level share a memoized result.

\subsection{Zero resolvents are not exceptions to the theory}

A resolvent may be the zero tensor, or a nonzero tensor may evaluate to zero
at the chosen pair of embeddings. Neither case invalidates Fourier inversion.
A zero summand contributes zero; it is never used as a denominator. The
implementation distinguishes formal zero from zero in the chosen embedding
when choosing branches. This is another advantage of using all the Fourier
components to reconstruct the requested element rather than insisting that one
particular component be a nonzero field generator.

\section{Branch selection using an exact annihilator}\label{sec:branches}

A root-of-unity multiplier is not optional decoration. Independent principal
$p$th roots of the $Q_k$ generally do not reconstruct $\beta$. The branch
choice must be tied to the concrete element whose field coordinates are known.

\subsection{The norm polynomial}

Represent $R_k(z)$ by $R(u,z)$ with rational coefficients reduced in $u$ modulo
$h$ and in $z$ modulo $\Phi_p$. Compute
\begin{equation}\label{eq:norm}
 P_R(V)=\Res_u\!\left(h(u),
       \Res_z\bigl(\Phi_p(z),V-R(u,z)\bigr)\right).
\end{equation}
It is a monic rational polynomial of degree $d(p-1)$ before repeated factors
are removed. It is the characteristic/norm polynomial of the multiplication
operator by $R$ in the finite algebra~\eqref{eq:tensor}. It annihilates every
evaluation of $R$ at a conjugate of $t$ and a primitive $p$th root of unity.
The code uses its monic square-free part.

\begin{lemma}\label{lem:allbranches}
For $1\le k<p$, every candidate
$\zeta_p^e Q_k^{1/p}$ in~\eqref{eq:branchresolvent} is a root of $P_R$.
\end{lemma}
\begin{proof}
If $R_k(\zeta_p)=0$, all candidates are zero. Otherwise the $p$ candidates
are precisely the $p$ values obtained by multiplying $R_k(\zeta_p)$ by a
$p$th root of unity. Equation~\eqref{eq:eigen} rotates that value by
$\zeta_p^k$ when $\sigma$ acts on the $L$ coordinates and fixes $z$.
Because $k$ is coprime to the prime $p$, its powers generate all these
multipliers. All those evaluations occur among the roots of the norm
polynomial. The equality $Q_k=R_k(\zeta_p)^p$ is already established by the
exact tensor power and the induction hypothesis for its coefficients.
\end{proof}

The proof is important: it justifies a common-annihilator root comparison for
\emph{both} expressions being compared. Passing an arbitrary numerical guess
to such a comparison would not be justified.

\subsection{Equality is easier when a polynomial is already known}

For a square-free rational polynomial $P$, distinct complex roots have an
effectively computable positive separation bound $\delta$. If two known roots
are evaluated with absolute error less than $\delta/8$, their approximate
distance is less than $\delta/4$ when they are equal and greater than
$3\delta/4$ when they are different. A threshold between these cases decides
equality. The evaluation tolerance comes from a theorem about $P$, not from a
user's expectation that 50 or 100 decimal places ought to be enough.

SymPy's \code{Poly.same\_root} is specifically an equality test for two known
roots \cite{sympy-poly}. In the pinned implementation it uses a Mignotte
separation bound and bounded-absolute-error evaluation. The package invokes it
only after a structural identity, a factor relation, or
Lemma~\ref{lem:allbranches} supplies the known-root premise. The degree-one
case is handled separately, because under that premise there is only one root.

This remains a trusted-CAS calculation, not a formally verified interval kernel
supplied with the package. In particular, the report does not claim that a
small independent proof assistant has checked SymPy's numerical evaluator.
The distinction from heuristic recognition is that the polynomial and required
separation are established algebraically before the comparison is made.

\subsection{Two levels of verification}

The general engine maintains construction invariants: exact field operations,
fixed-field membership, the resolvent-power identity, and a valid common
annihilator for each branch comparison. The resulting expression therefore
inherits its algebraic identity from the construction.

An independent verifier can instead forget this history, compute the minimal
polynomial $m_E$ of the final radical expression $E$, check $m_E\mid f$, and
then compare embeddings. The Python \code{verify} function does exactly that.
It is useful for regression tests, but can be dramatically slower: treating
correlated radicals independently may create a much larger temporary field.
The C5 experiment in Section~\ref{sec:tests} illustrates this distinction.

The Wolfram frontend performs an additional check
\begin{lstlisting}
RootReduce[answer - input] === 0
\end{lstlisting}
and returns no candidate unless it passes. A timed-out or unresolved check is
not silently replaced by numerical proximity.

\section{Correctness and finite completeness}\label{sec:complete}

\begin{theorem}\label{thm:complete}
Let $f\in\Q[x]$ be irreducible and let an embedding of one of its roots be
specified. With resource bounds removed and correct terminating exact
factorization, resultant, and algebraic-embedding primitives, the general
algorithm terminates. It produces a radical expression for the specified root
when that root is expressible by radicals, and otherwise proves
nonsolvability through its Galois group.
\end{theorem}
\begin{proof}
Section~\ref{sec:closure} constructs the splitting field in finitely many
adjunctions, with every primitive-element search terminating by
Lemma~\ref{lem:primitive}. Proposition~\ref{prop:autos} gives all automorphisms
by a finite search and exact polynomial tests. The finite multiplication table
then either detects a nontrivial perfect subgroup during the derived-subgroup
construction or gives the prime-index chain~\eqref{eq:chain}.

In the former case $G$ is not solvable, and the necessity direction of the
Galois solvability theorem excludes a radical expression for any root of this
irreducible polynomial. In the latter case use induction on fixed-field level.
The base field is $L^G=\Q$. For the inductive step,
Proposition~\ref{prop:descent} puts the trace and all coefficients of each
resolvent power in the preceding level. The induction hypothesis converts
them to radicals. There are finitely many possible branches, and
Lemma~\ref{lem:allbranches} supplies a correct annihilator for a terminating
algebraic equality test that selects the actual one. Fourier inversion then
proves the result equals the original field element. The induction constructs
all roots, and an exact comparison selects the requested embedding.
\end{proof}

\begin{corollary}
A resource cutoff during this calculation proves nothing about radical
solvability. In particular, \code{ResourceLimit} and \code{NotSolvable} are
logically different outcomes.
\end{corollary}

The theorem does not require all extensions in the raw root-adjunction tower
to be radical extensions. They usually are not presented that way. That tower
is a device for computing the normal field and its action. The radical tower
is built only afterward from its solvable group.

Nor does the theorem assert that every intermediate field has a primitive
generator of lower absolute degree than the original input. Normal closures
and cyclotomic coordinates can be larger than the original root field. Trying
to force monotone decrease of the \emph{absolute minimal-polynomial degree} at
every discovery step is not a valid general completeness principle. What
strictly decreases during radical synthesis is the subgroup-chain level.

\section{Implementation details that affect correctness}\label{sec:implementation}

\subsection{Rational vectors and orientation}

The reference implementation uses SymPy's exact algebraic-number polynomial
representation. Its low-level \code{ANP.to\_list()} coefficients are in
descending order, while tensor tuples are in ascending powers of $z$. The
conversion helpers isolate that difference. The defining minimal polynomial
is explicitly supplied to \code{QQ.algebraic\_field}; it is not recomputed from
an already complex nested expression.

Automorphism composition uses polynomial substitution, not arithmetic on
approximate roots. Group elements are integers indexing the multiplication
table. Subgroups are sets of these indices. Memoization keys for field elements
are their exact coefficient tuples and their descent level; different
conjugates are not conflated merely because their minimal polynomials agree.

The code accesses some SymPy representation-level interfaces, including
\code{Poly.rep} and algebraic-number polynomial lists. The dependency is pinned
to SymPy 1.14.0, the version used for testing. Porting to another representation
or library should preserve these orientation and normalization invariants.

\subsection{No arbitrary tensor division}

Tensor multiplication forms the coefficient convolution and reduces terms
from highest degree downward with $z^{p-1}=-\sum_{j=0}^{p-2}z^j$.
Exponentiation uses repeated squaring. At no point does the code ask for the
inverse of a nonzero tensor element. Field-element inverses, such as those
needed for exact linear factors, occur only in the genuine field $L$.
A regression test explicitly multiplies two nonzero zero divisors in the
$p=3$ tensor algebra and checks that the product is zero.

\subsection{Expression DAGs, not executable source strings}

The Python backend serializes output as a directed acyclic graph. Its node
tags are \code{Q}, \code{I}, \code{Add}, \code{Mul}, and \code{Pow}. Rationals
carry signed decimal numerator and denominator strings; references are
zero-based indices of earlier nodes. The root index points to the final value.
Every power exponent must decode to a rational number.

The Wolfram decoder validates tags, arities, rational denominators, backward
references, and the final radical grammar. It uses integer parsing and
arithmetic constructors, not \code{ToExpression}. The Python CLI similarly
accepts rational coefficient data rather than expressions for \code{eval}.
This also avoids copying repeated subexpressions into the interchange format.
It does not guarantee that every later CAS pretty-printer will preserve
sharing when displaying the full formula.

\subsection{Root indices are not portable}

The Python API defines \code{index} by SymPy's \code{CRootOf(p,index)} ordering,
including multiplicities when a selected root is requested. Its
\code{all\_roots} API returns distinct roots with no promised ordering.

The Wolfram frontend never translates a Mathematica root index into a Python
root index. It sends the minimal polynomial, requests all its roots, and selects
a candidate by exact equality to the original Wolfram expression. The extra
work buys a simple and reliable cross-system embedding contract. A more
optimized future protocol could transmit an exact isolating rectangle instead,
but approximate coordinates alone would not justify a change in this contract.

\subsection{Failure and time-limit semantics}

The important statuses are shown in Table~\ref{tab:status}. Python exposes
status-bearing exception classes; the CLI serializes analogous strings; the
Wolfram frontend uses \code{Failure} expressions.

\begin{table}[ht]
\centering\small
\begin{tabularx}{\linewidth}{@{}lX@{}}
\toprule
Status & Meaning \\
\midrule
\code{Success} & A radical expression was constructed with the method's checks. \\
\code{NotSolvable} & A nonsolvable Galois group was established for the relevant irreducible factor. \\
\code{ResourceLimit} & A field-degree budget or computation deadline was exceeded. \\
\code{InvalidInput} & Input was not in the exact rational-polynomial contract, or options/data were malformed. \\
\code{Unresolved} & The explicitly requested native-only structural search did not yield a verified result. \\
\code{BackendError} & An underlying computation failed; no mathematical impossibility follows. \\
\code{VerificationFailed} & The frontend did not establish equality to the requested embedding. \\
\bottomrule
\end{tabularx}
\caption{Operational outcomes must not be collapsed into nonsolvability.}\label{tab:status}
\end{table}

The in-process Python \code{max\_seconds} budget is cooperative: it is checked
between operations, so a single factorization or root-isolation call can run
past the deadline. The bundled \code{python/cli.py} additionally supports
\code{hard\_timeout\_seconds}; it starts an isolated worker and terminates that
worker when the deadline expires. The Wolfram frontend uses that supervisor
when its \code{TimeConstraint} is finite.

There is still no single global Wolfram deadline: native candidate generation,
the Python worker, and final per-candidate exact verification have separate
stages. The last stage is controlled by
\code{"VerificationTimeConstraint"}; its default is \code{Infinity}. These semantics are intentional and documented rather than
hidden behind a misleading total-time promise.

\section{Using the package}\label{sec:usage}

\subsection{Wolfram Language}

Keep the unpacked directory layout intact and load the file:
\begin{lstlisting}
Get["C:\\path\\to\\RadicalSolve\\wolfram\\RadicalSolve.wl"];
r = Root[-1-#^2-#^3+#^4+#^6&, 2];
a = Radicalize[r, Method -> "Native"];
VerifyRadical[r, a]
RadicalReport[r, Method -> "Native"]
\end{lstlisting}
The native route includes the structural methods and low-degree
\code{ToRadicals} formulas. Both original examples have native structural
routes, so this use does not require Python.

For the general route, install the pinned Python dependency using the chosen
interpreter, then supply its executable name or absolute path:
\begin{lstlisting}
RadicalReport[Root[#^3 - 3# + 1&, 1],
  Method -> "Galois",
  "PythonExecutable" -> "python",
  "MaxFieldDegree" -> 96,
  TimeConstraint -> 120,
  "VerificationTimeConstraint" -> 120]
\end{lstlisting}
\code{Method -> Automatic} tries native methods first and then the Python
backend. \code{Method -> "Native"} never launches it.
\code{Method -> "Galois"} bypasses structural reduction for non-radical input.
An expression already in the radical grammar is simply returned.

\code{RadicalReport} returns an association containing the expression, method,
and verification result. Backend construction metadata is included when
applicable. The default field-degree cap is 96; use \code{Infinity} to remove it.

\subsection{Python}

From the archive root, install or use the unpacked module:
\begin{lstlisting}
python -m pip install -e .
python -m unittest discover -s tests -v
\end{lstlisting}
Alternatively add the \code{python/} directory to \code{PYTHONPATH}; no package
installation is required for the bundled CLI.

\begin{lstlisting}
from sympy import symbols, Poly
from radicalsolve import Solver, radicalize, verify
x = symbols("x")
p = Poly(x**6+x**4-x**3-x**2-1, x)
r = radicalize(p, index=1)
assert verify(p, r.expression, index=1)
print(r.expression)
print(r.method)

session = Solver(method="galois", max_field_degree=96)
q = Poly(x**5+x**4-4*x**3-3*x**2+3*x+1, x)
roots = session.all_roots(q)
print(session.events)
\end{lstlisting}
\code{max\_field\_degree=None} removes the Python degree cap.
\code{Solver.root} first selects the irreducible factor containing the requested
root. This is why it can return a solvable root of
$(x^2-2)(x^5-x-1)$ without being blocked by the nonsolvable quintic factor.
In contrast, \code{all\_roots} requires all irreducible factors to be solvable.

\code{Result.expression} is the result, \code{Result.method} is the route,
and \code{Result.record} is an audit record. The record is \emph{not} a
standalone, independently replayable proof certificate. This version performs
its construction checks during execution and offers independent minimal-
polynomial verification separately. A compact portable certificate checker is
an explicit future engineering improvement, not an already supplied feature.

\section{Executed tests and observed limitations}\label{sec:tests}

\subsection{Environment and reproducibility}

The reference code was executed with Python 3.13 and SymPy 1.14.0. The archive
contains the actual \code{unittest} transcript and a separate general-$S_4$
stress transcript. The declared Python requirement is 3.10 or later; only the
stated environment was exercised here.

No usable Wolfram kernel was available. Both attempted Wolfram connector calls
failed before evaluation. Consequently the supplied \code{.wl} and
\code{.wlt} files have source-level review but \emph{no claimed Wolfram runtime
test}. The Python engine is the executed reference, and the native formulas
are also justified mathematically in this report.

\subsection{Coverage}

The default suite exercises the two original selected roots with independent
minimal-polynomial and embedding checks; it also matches every conjugate from
the two structural constructions. Additional tests cover translated inputs,
power compositions, rational roots, multiplicities, a mixed reducible
solvability example, wrong conjugates, invalid coefficients, invalid indices,
tensor zero divisors, data encoding, and process deadlines.

The forced-general-method families are listed in Table~\ref{tab:tests}.
These tests bypass the fast structural layer; otherwise, examples such as
$x^4-2$ would say nothing about the general engine.

\begin{table}[ht]
\centering\small
\begin{tabularx}{\linewidth}{@{}lXcl@{}}
\toprule
Group & Polynomial & $[L:\Q]$ & Independent final check \\
\midrule
$C_2$ & $x^2+1$ & 2 & minimal polynomial \\
$C_3$ & $x^3-3x+1$ & 3 & minimal polynomial \\
$S_3$ & $x^3-2$ & 6 & minimal polynomial \\
$V_4$ & $x^4-10x^2+1$ & 4 & minimal polynomial \\
$C_4$ & $x^4+x^3+x^2+x+1$ & 4 & minimal polynomial \\
$D_4$ & $x^4-2$ & 8 & minimal polynomial \\
$C_5$ & $x^5+x^4-4x^3-3x^2+3x+1$ & 5 & numerical residual only \\
\bottomrule
\end{tabularx}
\caption{Executed general-descent families. All rows also execute the internal
exact algebraic invariants and separation-based branch checks. ``Numerical
residual only'' describes the additional, independent final check, not the
construction algorithm. $D_4$ denotes the group of order eight.}\label{tab:tests}
\end{table}

The $C_5$ example is the real cyclotomic degree-five polynomial associated with
order eleven. It can be derived by dividing $1+z+\cdots+z^{10}$ by $z^5$ and
expressing the symmetric terms in $x=z+z^{-1}$. This gives a useful independent
mathematical description of its solvability, although the test forces the
implemented Fourier route rather than using that description as a shortcut.

The default test run reports the actual pass/skip count in the included
transcript. The optional forced-$S_4$ integration test is skipped by default.
The ordinary automatic/classical route for a generic quartic is nevertheless
included and passes its tests.

\subsection{What the stress run did and did not establish}

For $x^4-x-1$, a forced-general-method exploratory run constructed the degree-24
splitting field, recovered its automorphisms from weights $(1,2,3)$, and obtained
the chain of subgroup orders
\begin{equation}
 24,12,4,2,1.
\end{equation}
The corresponding prime indices were $2,3,2,2$. That run did not complete
radical output within its cooperative time budget. A sampled stack was inside
complex root isolation used by a branch comparison. It is recorded as a
resource-limited run, not as a passing full integration test.

For the cyclic quintic, constructing all five radical expressions in an
exploratory run took about a second in this environment. A separate attempt
to reconstruct the minimal polynomial of the expanded result ran much longer
and was stopped. The completed regression test therefore uses the construction
checks and an additional high-precision residual rather than claiming that the
independent reconstruction finished. This is an observed verification-cost
issue, not a claim of a mathematical ambiguity in the formula.

The large inverse-beta-related notebook examples, including degrees beyond
those in Table~\ref{tab:tests}, are retained as source material. They were not
silently promoted into a ``passed'' benchmark set. A useful next benchmark
campaign would compare normal-closure degree, coefficient growth, descent size,
and verification time on those exact inputs.

\section{Complexity and further improvements}\label{sec:improvements}

\subsection{The main costs}

Write $d=[L:\Q]=|G|$. The input degree $n$ alone is a poor predictor of cost:
$d$ can approach $n!$. A full multiplication table requires $O(d^2)$ entries,
and constructing it involves polynomial substitutions in a degree-$d$ field.
The primitive-image search examines at most $n!/(n-r)!$ ordered tuples when
$r$ roots occur in the primitive sum.

At a cyclic step of prime index $p$, tensor coordinates have dimension $p-1$
over $L$. Their rational norm polynomial has degree $d(p-1)$ before square-free
reduction. In a naive recursion, a level can create up to $(p-1)^2$ coefficient
subproblems plus a trace subproblem. Memoization helps substantially but does
not turn this into a general polynomial-time output algorithm. Rational
coefficient heights and algebraic-number isolation costs are equally important.

The resulting formula need not be close to minimal. A short radical can acquire
a large expression because the construction used a complicated primitive
coordinate system. The structural layer exists precisely to avoid that
unnecessary detour whenever an exact cheap identity is available.

\subsection{Improvements that preserve the proof structure}

The most useful extensions are computational, not changes to the solvability
criterion. Better number-field factorization and root-isolation backends would
address the observed bottlenecks. Smaller primitive elements and reduced bases
could limit coefficient growth. Acting with subgroup generators rather than
checking all subgroup elements could reduce repeated fixed-field tests after
an exact generating set has been certified.

A more compact descent can choose a suitable nonzero eigenvector and a relative
power basis, then express several requested elements in that basis. That may
share radicals more effectively than reconstructing each element from all
Fourier components. It introduces extra nonvanishing and basis-independence
obligations, especially when a cyclotomic tensor splits. The current version
accepts the larger formulas in exchange for a simple division-free argument.

The native structural recognizers could also be broadened to rational-function
decompositions, additional resolvent families, or systematically certified
subfield searches inspired by the notebook. Those methods should remain
optimizations: a failed structural search must not overwrite the general
solvability status.

Finally, a portable proof record should include the field polynomial,
automorphism images, subgroup chain, resolvent-power identities, and selected
branch witnesses. An independent checker could then replay exact identities
without rediscovering a large minimal polynomial from a flat formula. The
present audit metadata is not yet that full object.

\section{Conclusion}

A systematic converter needs more than a longer list of radical formulas. It
needs the normal field and its actual action, a descent invariant that truly
decreases, and an embedding contract that survives every radical extraction.
The supplied implementation gives those pieces concrete representations.

The two motivating examples receive short structural answers. The general
engine implements a finite solvable-Galois construction, including prime cyclic
steps beyond cubic and quartic formulas. The tensor formulation avoids assuming
that roots of unity are disjoint from the splitting field, and the norm argument
makes branch comparisons legitimate before any numerical evaluation is used.
The delivered code is a tested reference implementation with explicit resource
and validation boundaries, rather than a claim that every large solvable input
will already have an economical answer.

\clearpage
\appendix
\section{Compact algorithm descriptions}

\subsection{Normal closure and automorphisms}
\begin{lstlisting}
h := irreducible input polynomial f
weights := [1]
repeat:
    K := Q[t]/h
    factor f over K
    if all factors are linear: break
    choose a nonlinear irreducible factor g of degree e
    for c = 1,2,...:
        H := Res_u(h(u), c^e g((T-u)/c,u))
        if degree(H) = e degree(h) and H is square-free:
            h := monic(H)
            append c to weights
            break

roots := linear roots of f in K
images := distinct weighted sums of distinct roots satisfying h(image)=0
require number(images) = degree(h)
construct automorphism multiplication by exact substitution
\end{lstlisting}

\subsection{One radical-conversion call}
\begin{lstlisting}
Rad(beta, level):
    require beta fixed by H[level]
    if beta is rational: return beta
    if beta fixed by H[level-1]: return Rad(beta, level-1)
    p, sigma := current prime quotient and generator
    pieces := [Rad(trace(beta), level-1)]
    for k = 1,...,p-1:
        R := sum_j z^(-k j) sigma^j(beta) modulo Phi_p(z)
        if R is formally zero: append 0; continue
        P := square-free rational norm polynomial of R
        handle a linear P or zero in the selected embedding directly
        q := R^p, computed without division in the tensor algebra
        require every coefficient of q fixed by H[level-1]
        Q := evaluate recursively converted q at z = zeta_p
        select e with R(zeta_p) = zeta_p^e Q^(1/p)
            using the common polynomial P and root separation
        append zeta_p^e Q^(1/p)
    return sum(pieces)/p
\end{lstlisting}

\section{Files and rebuilding}

The archive's top-level \code{README.md} gives installation and command examples.
The implementation files are under \code{python/radicalsolve/}; the direct
supervised entry point is \code{python/cli.py}. The native Wolfram frontend is
\code{wolfram/RadicalSolve.wl}. Tests and actual transcripts are under
\code{tests/}. The unchanged original notebook and its source hashes are under
\code{provenance/}.

The article can be rebuilt in its directory with
\begin{lstlisting}
latexmk -pdf radicalsolve.tex
\end{lstlisting}
It uses ordinary LaTeX packages and Latin Modern fonts. No external images,
commercial fonts, or downloaded font files are required. New implementation
and report material are provided under MIT-0; the original notebook and
third-party dependencies retain their existing rights.

\begin{thebibliography}{9}
\bibitem{question}
Vladimir Reshetnikov.
\emph{Why is ToRadicals not able to handle all cases? Is there a workaround?}
Mathematica Stack Exchange, question 34011 and associated discussion.
\url{https://mathematica.stackexchange.com/questions/34011/}.
Accessed 8 September 2026.

\bibitem{notebook}
Vladimir Reshetnikov.
\emph{FindExtension.nb}, supplied in \code{FindExtension.zip}, received
8 September 2026. Unchanged source and SHA-256 hashes included in the archive.

\bibitem{milne}
J. S. Milne. \emph{Fields and Galois Theory}, version 5.10, 2022.
Especially the cyclic-extension and solvability material in Chapter 5.
\url{https://www.jmilne.org/math/CourseNotes/FT.pdf}.

\bibitem{li}
Song Li. \emph{An Algorithm for Solving Solvable Polynomial Equations of
Arbitrary Degree by Radicals}. 2022. arXiv:2203.11602.
\url{https://arxiv.org/abs/2203.11602}.

\bibitem{wolfram-toradicals}
Wolfram Research. \emph{ToRadicals}, Wolfram Language documentation.
\url{https://reference.wolfram.com/language/ref/ToRadicals.html}.
Accessed 8 September 2026.

\bibitem{wolfram-exact}
Wolfram Research. \emph{ExactNumberQ}, Wolfram Language documentation.
\url{https://reference.wolfram.com/language/ref/ExactNumberQ.html}.
Accessed 8 September 2026.

\bibitem{sympy-fields}
SymPy Development Team. \emph{Number Fields}, SymPy 1.14.0 documentation.
\url{https://docs.sympy.org/latest/modules/polys/numberfields.html}.
Accessed 8 September 2026.

\bibitem{sympy-poly}
SymPy Development Team. \emph{Polynomials Manipulation Module Reference},
including \code{Poly.same\_root}, SymPy 1.14.0 documentation and implementation.
\url{https://docs.sympy.org/latest/modules/polys/reference.html}.
Accessed 8 September 2026.
\end{thebibliography}
\end{document}
